\chapter{\ifenglish Conclusions and Discussions\else บทสรุปและข้อเสนอแนะ\fi}

\section{\ifenglish Conclusions\else สรุปผล\fi}

ปริญญานิพนธ์ฉบับนี้ได้นำเสนอการพัฒนา "ระบบจัดการห่วงโซ่อุปทานด้วยบล็อกเชนและไอโอที " โดยมีวัตถุประสงค์เพื่อยกระดับความปลอดภัยและความน่าเชื่อถือของข้อมูลอุณหภูมิในกระบวนการขนส่ง จากผลการดำเนินงานสามารถสรุปได้ว่า ระบบที่พัฒนาขึ้นสามารถทำงานได้บรรลุตามวัตถุประสงค์ทุกประการ 

ระบบสามารถอ่านค่าอุณหภูมิผ่านอุปกรณ์ตรวจวัดและส่งข้อมูลตามเวลาจริงผ่านโปรโตคอลการสื่อสาร โดยมีความหน่วงเวลาเฉลี่ยเพียง 1.2 วินาที ข้อมูลทั้งหมดถูกบันทึกลงในเครือข่ายบล็อกเชนแบบอนุญาต ซึ่งรับประกันความสมบูรณ์ของข้อมูลได้ร้อยละ 100 โดยไม่มีการสูญหายหรือถูกดัดแปลงแก้ไข นอกจากนี้หน้ากระดานสรุปข้อมูลยังสามารถแสดงผลข้อมูลย้อนหลังและแจ้งเตือนเมื่ออุณหภูมิออกนอกเกณฑ์ที่กำหนดได้อย่างแม่นยำ

\textbf{ข้อจำกัดของระบบ} \\
แม้ระบบจะสามารถทำงานได้ตามเกณฑ์มาตรฐาน แต่ยังมีข้อจำกัดบางประการที่พบจากการพัฒนาและทดสอบ ดังนี้
\begin{enumerate}
    \item \textbf{ข้อจำกัดด้านเครือข่ายของอุปกรณ์ฮาร์ดแวร์:} อุปกรณ์ไอโอที ที่ใช้ในโครงงานนี้ยังคงพึ่งพาการเชื่อมต่อผ่านระบบเครือข่ายไร้สาย เป็นหลัก ทำให้ไม่เหมาะกับการนำไปติดตั้งบนยานพาหนะขนส่งที่ต้องเคลื่อนที่ตลอดเวลาในพื้นที่ที่ไม่มีสัญญาณ
    \item \textbf{ข้อจำกัดด้านโครงสร้างพื้นฐานของบล็อกเชน:} ในช่วงการพัฒนาและทดสอบนี้ เครือข่ายบล็อกเชนถูกจำลองและติดตั้งรวมอยู่บนเครื่องแม่ข่ายคลาวด์เพียงโหนดเดียว แม้ในเชิงตรรกะจะมีการแยกองค์กร แต่ในเชิงกายภาพยังถือว่ามีโครงสร้างพื้นฐานแบบรวมศูนย์ ระดับหนึ่ง ซึ่งในสภาพแวดล้อมจริงควรแยกโหนดไปตามองค์กรต่าง ๆ
    \item \textbf{ข้อจำกัดด้านการแจ้งเตือน:} ระบบแจ้งเตือนในปัจจุบันเป็นการทำงานแบบตรวจสอบตามเงื่อนไขที่กำหนด เมื่ออุณหภูมิเกินค่าที่กำหนดจึงจะทำการแจ้งเตือน แต่ยังไม่มีระบบการวิเคราะห์เชิงคาดการณ์ล่วงหน้า
\end{enumerate}

\section{\ifenglish Challenges\else ปัญหาที่พบและแนวทางการแก้ไข\fi}

ในการทำโครงงานนี้ พบว่าเกิดปัญหาหลัก ๆ และได้ดำเนินการแก้ไข ดังนี้

\begin{enumerate}
    \item \textbf{ปัญหาข้อมูลตกค้างในระบบจัดการคอนเทนเนอร์}
    \begin{itemize}
        \item ปัญหาที่พบ: ในระหว่างการพัฒนาและเริ่มต้นเครือข่ายบล็อกเชนใหม่ มักพบข้อผิดพลาดทำให้ไม่สามารถสร้างช่องทางการสื่อสาร หรือบันทึกข้อมูลได้ เนื่องจากมีข้อมูลบัญชีธุรกรรมเดิมตกค้างอยู่ในระบบ
        \item แนวทางการแก้ไข: พัฒนาชุดคำสั่งสำหรับการล้างข้อมูลระดับลึก โดยบังคับหยุดและลบคอนเทนเนอร์ เครือข่าย และข้อมูลที่ตกค้าง รวมถึงลบโฟลเดอร์เก็บข้อมูลประจำตัวทิ้งทั้งหมด เพื่อให้ระบบสามารถสร้างข้อมูลการระบุตัวตนและบล็อกเริ่มต้นใหม่ได้อย่างสมบูรณ์
    \end{itemize}
    \item \textbf{ปัญหาการเข้าถึงไฟล์ใบรับรองอิเล็กทรอนิกส์}
    \begin{itemize}
        \item ปัญหาที่พบ: เมื่อมีการตั้งค่าเครือข่ายบล็อกเชนใหม่ ชื่อไฟล์ของใบรับรองอิเล็กทรอนิกส์ และกุญแจส่วนตัวสำหรับการทำธุรกรรมมักจะถูกสุ่มสร้างใหม่ ทำให้ชุดคำสั่ง Node.js เดิมที่ระบุชื่อไฟล์แบบตายตัวไม่สามารถค้นหาไฟล์พบ
        \item แนวทางการแก้ไข: ปรับปรุงชุดคำสั่งในส่วนของซอฟต์แวร์ชั้นกลาง โดยเขียนฟังก์ชันการค้นหาอัตโนมัติ เพื่อให้ระบบเข้าไปอ่านไฟล์แรกสุดที่พบในโฟลเดอร์เก็บกุญแจและใบรับรองโดยอัตโนมัติ ทำให้ระบบเชื่อมต่อข้อมูลประจำตัวได้เสมอโดยไม่ต้องแก้ไขชุดคำสั่งใหม่ทุกครั้ง
    \end{itemize}
    \item \textbf{ปัญหาความเสถียรของการเชื่อมต่ออินเทอร์เน็ต}
    \begin{itemize}
        \item ปัญหาที่พบ: หากสัญญาณเครือข่ายไร้สายขัดข้องหรือขาดหายชั่วคราว อุปกรณ์ตรวจวัดจะไม่สามารถส่งข้อมูลไปยังตัวกลางรับส่งข้อความได้
        \item แนวทางการแก้ไข: ได้ออกแบบระบบให้รองรับการกู้คืนอัตโนมัติ โดยเมื่อสัญญาณกลับมาเป็นปกติ ระบบจะสามารถเชื่อมต่อและกลับมาทำงานต่อได้เองภายใน 60 วินาที โดยไม่จำเป็นต้องให้ผู้ใช้งานทำการเริ่มต้นระบบใหม่ด้วยตนเอง
    \end{itemize}
\end{enumerate}

\section{\ifenglish Suggestions and further improvements\else ข้อเสนอแนะและแนวทางการพัฒนาต่อ\fi}

ข้อเสนอแนะเพื่อพัฒนาโครงงานนี้ต่อไป มีดังนี้

\begin{enumerate}
    \item \textbf{การปรับปรุงส่วนติดต่อผู้ใช้งานและประสบการณ์ผู้ใช้งาน:} อ้างอิงจากผลการประเมินของผู้ใช้งาน ควรเพิ่มการรองรับภาษาไทยควบคู่กับภาษาอังกฤษในหน้ากระดานสรุปข้อมูล ปรับขนาดของกล่องแสดงผลข้อมูลอุณหภูมิให้มีขนาดใหญ่ขึ้นเพื่อให้มองเห็นได้ชัดเจน และปรับพื้นที่การจัดวางให้ดูทันสมัยและมีเอกลักษณ์มากยิ่งขึ้น
    \item \textbf{การพัฒนาด้านฮาร์ดแวร์เพื่อการติดตามแบบเคลื่อนที่:} ควรเปลี่ยนไปใช้โมดูลการสื่อสารผ่านเครือข่ายโทรศัพท์เคลื่อนที่ เช่น NB-IoT หรือ 4G/5G แทนการใช้เครือข่ายไร้สาย รวมทั้งเพิ่มอุปกรณ์ระบุตำแหน่ง เพื่อให้สามารถติดตามข้อมูลอุณหภูมิควบคู่ไปกับตำแหน่งพิกัดของรถขนส่งได้ตามเวลาจริง
    \item \textbf{การประยุกต์ใช้ปัญญาประดิษฐ์:} ควรนำข้อมูลประวัติอุณหภูมิที่ถูกบันทึกไว้อย่างปลอดภัยในบล็อกเชน มาใช้เป็นชุดข้อมูลสำหรับฝึกสอนโมเดลการเรียนรู้ของเครื่อง เพื่อพัฒนาระบบการวิเคราะห์เชิงคาดการณ์ ซึ่งจะช่วยแจ้งเตือนก่อนที่อุณหภูมิจะสูงถึงจุดวิกฤต
    \item \textbf{การขยายเครือข่ายบล็อกเชนในระดับกลุ่มองค์กร:} เพื่อให้เกิดการกระจายศูนย์ อย่างแท้จริง ควรพัฒนาโครงงานต่อโดยการจัดเตรียมเครื่องแม่ข่ายคลาวด์แยกกันสำหรับแต่ละองค์กรในห่วงโซ่อุปทาน (เช่น โหนดของฟาร์ม โหนดของบริษัทขนส่ง โหนดของร้านค้าปลีก) และให้แต่ละโหนดทำการตรวจสอบและรับรองธุรกรรมร่วมกัน
    \item \textbf{การรองรับการขยายขนาดเครือข่ายบล็อกเชน:} ในอนาคตหากมีองค์กรเข้าร่วมในห่วงโซ่อุปทานเพิ่มมากขึ้น ควรพิจารณาปรับเปลี่ยนสถาปัตยกรรมการจัดการคอนเทนเนอร์ จากเดิมที่ปฏิบัติงานบนเครื่องแม่ข่ายเดี่ยว ไปสู่การใช้ระบบบริหารจัดการคอนเทนเนอร์ระดับอุตสาหกรรม เช่น คูเบอร์เนทีส เพื่อให้สามารถขยายหรือลดขีดความสามารถของโหนดในเครือข่ายบล็อกเชนได้อย่างยืดหยุ่น และสามารถรองรับการทำงานข้ามผู้ให้บริการระบบคลาวด์แบบพหุคูณได้อย่างมีประสิทธิภาพ
    \item \textbf{การจัดการการนำเข้าข้อมูลปริมาณมหาศาลจากอุปกรณ์ไอโอที:} เมื่อจำนวนอุปกรณ์ตรวจวัดเพิ่มขึ้นสู่ระดับอุตสาหกรรม ควรมีการยกระดับระบบตัวกลางรับส่งข้อความ ไปสู่บริการระดับองค์กรที่รองรับสภาพพร้อมใช้งานสูง รวมถึงการจัดสมดุลภาระงาน ที่ฝั่งซอฟต์แวร์ชั้นกลาง เพื่อกระจายภาระการประมวลผลและป้องกันปัญหาคอขวด ของข้อมูลก่อนนำส่งเข้าสู่ระบบบล็อกเชน
    \item \textbf{การบริหารจัดการข้อมูลนอกเครือข่ายบล็อกเชน:} เพื่อป้องกันมิให้ขนาดของบัญชีธุรกรรมขยายตัวเกินความจำเป็นเมื่อเวลาผ่านไป ควรปรับสถาปัตยกรรมไปสู่ระบบจัดเก็บข้อมูลแบบผสมผสาน โดยนำข้อมูลอุณหภูมิที่มีขนาดใหญ่และข้อมูลรายละเอียดเชิงลึกไปจัดเก็บไว้ในฐานข้อมูลนอกเครือข่ายบล็อกเชน และบันทึกเฉพาะค่าแฮชของข้อมูลเหล่านั้นลงบนเครือข่ายบล็อกเชน ซึ่งวิธีการนี้จะช่วยเพิ่มความเร็วในการสืบค้นข้อมูลย้อนหลัง โดยยังคงรักษาคุณสมบัติการป้องกันการดัดแปลงแก้ไขข้อมูลไว้ได้อย่างครบถ้วน
\end{enumerate}